\documentclass{article}
\usepackage[preprint]{neurips_2026}
\renewcommand{\ttdefault}{lmtt} % Type 1 typewriter; keeps the template's Times
\usepackage[T1]{fontenc}
\usepackage[utf8]{inputenc}
\usepackage{booktabs}
\usepackage{amsmath}
\usepackage{graphicx}
\usepackage{xcolor}
\usepackage{microtype}
\usepackage{url}
\newcommand{\repourl}{\url{https://github.com/jacksonmlukas/connections-rl}}

\title{Reliable and Invalid: Anatomy of a Gamed Training Task}
\workshoptitle{}
\author{%
  Jackson Lukas \\
  \texttt{jlukas3313@gmail.com} \\
}

\begin{document}
\maketitle

\begin{abstract}
Verifiable rewards are trusted because they are deterministic: a program
either awards the points or it does not. We post-train Qwen2.5-7B-Instruct
with GRPO on NYT Connections under a reward with a cheap structural component
(well-formed output) and an expensive semantic component (correct
groupings). On a strictly chronological test set of 162 puzzles (648 group
slots), scored in one serving session, capability runs $26 \rightarrow 52
\rightarrow 80 \rightarrow 4$ slots across the untrained Instruct model,
supervised fine-tuning, GRPO step 50, and the GRPO endpoint: paired stage
gains of $+0.161$ and $+0.173$ on the
0--4 scale, then a fall to within $3\times$ of the 1.4-slot chance floor of a uniformly drawn valid partition, $-0.136$ $[-0.198, -0.080]$ below the untrained model,
reproduced across three seeds. The \emph{training} reward reads the
destruction as mastery: it reaches its 1.6 ceiling by step 125 with zero
across-rollout variance and stays there. Scoring the endpoint offline on
its own training puzzles --- 21 of 3{,}228 slots --- rules out memorization
and exposes the cause: a one-line bug in our GRPO data loader presented
each board's sixteen words in answer-key order, so the training task, as
presented, was solvable by copying the words four at a time. The collapsed policy does exactly that --- 91.4\% of the groups it emits on training boards, and 92.4\% on test boards, are consecutive quadruples of its prompt; a copy of a shuffled board is a random valid partition, which is why the endpoint lands near the chance floor. Entropy collapse, cross-seed
convergence, and the reward ceiling are signatures of one deterministic
copy rule; the structural climb is older, over half of it predating any
measurable copying. The reward function was correct
on every string it scored; the task presentation was gameable, and every
in-sample signal read the gaming as mastery. The same reward held out, and
three blind LLM judges, catch the reversal; four lines of string matching
then identify the rule. Determinism certifies the scorer, not the
construct: hold the reward out.
\end{abstract}

\section{Introduction}
\label{sec:intro}

Post-training against verifiable rewards is attractive for a reason that
sounds decisive: the reward is a program. It returns exactly what it was
written to return, and the argument that this
buys immunity from reward hacking\footnote{We follow the common usage
\emph{reward hacking} throughout; \citet{skalse2022}, whose formalization we
adopt, publish it under the name \emph{reward gaming}.} rests on that
reliability. This paper is about the gap reliability does not close: the gap
between the reward and the construct it stands in for. We train on NYT
Connections --- sixteen words, partitioned into four groups of four ---
chosen because well-formedness and correctness are unusually separable: a
response can be a perfectly valid partition and completely wrong. Our
reward pays for both; two of its five components can be checked by reading
the string, three require the answer key.

For fifty steps it works. In one serving session, held-out groups
recovered run $26 \rightarrow 52 \rightarrow 80$ across the untrained
Instruct model, the supervised checkpoint GRPO is initialized from, and
step 50 --- paired stage gains of $+0.161$ $[+0.056, +0.272]$ and $+0.173$
$[+0.074, +0.284]$, the reinforcement-learning stage contributing as much
as supervised fine-tuning (\S\ref{sec:control}). Over the following
353 steps the same run falls to 4 slots, below the untrained model and onto
the 1.4-slot floor of a uniformly drawn valid partition, while
the invalid-output rate falls by an order of magnitude and the training
reward reaches its 1.6 maximum and averages 1.5814 thereafter, with
zero variance across the $K=8$ rollouts on $\sim$95\% of prompts. A
practitioner watching that curve saw a solved task. The curve was not lying
about the reward; it was lying about the task. A control run scoring the
endpoint on its own training puzzles returned 21 of 3{,}228 slots --- not
memorization --- and led to the mechanism: our GRPO data loader, unlike our
SFT and evaluation code, presented each board's words in answer-key order,
and the converged policy simply copies the prompt four words at a time
(91--92\% of emitted groups; \S\ref{sec:control}). The failure is not in the reward
function, which is correct; it is in what the training distribution let
the reward certify, and in the practice of reading the certificate in
sample.

\paragraph{Contributions.} A fully verifiable case study in which GRPO
improves and then destroys the target construct, ending below the untrained
model across three seeds (\S\ref{sec:arc}); a complete post-hoc diagnosis:
an answer-ordered prompt bug made the training task gameable, and the
converged policy is a prompt-order copier on 91--92\% of its emitted groups
(\S\ref{sec:control}); an account in which entropy collapse, the reward
ceiling, and cross-seed convergence are signatures of that one rule
(\S\ref{sec:mechanism}), surviving a
normalization ablation (\S\ref{sec:ablation}); and the lesson demonstrated
end to end: every in-sample signal read gaming as mastery; the reward held
out, and three blind judges, caught it
(\S\ref{sec:control}--\S\ref{sec:judge}).

\section{Task, pipeline, and the reward's asymmetry}
\label{sec:setup}

\paragraph{Task and data.} 1{,}078 dated NYT Connections puzzles (1{,}077
after validation), split strictly chronologically: 807 train (2023-06-12 to
2025-08-27), 108 validation, 162 test (2025-12-15 to 2026-05-29); every
test puzzle postdates every training puzzle. Date regressions of
groups-correct across the test window are flat for both base ($-0.158$
groups/year $[-0.634, +0.292]$) and the trained policy ($+0.035$ $[-0.271,
+0.329]$): contamination controlled, though the intervals are wide and only
large drift would show.

\paragraph{Pipeline.} Qwen2.5-7B-Instruct (and 1.5B for the scale contrast),
rank-16 LoRA. SFT for 3 epochs on the training split, then GRPO
\citep{deepseekmath} from the SFT checkpoint: $K=8$ sampled completions per
puzzle at temperature 0.9, group-relative advantages, $\beta=0.001$ against
the SFT reference, learning rate $5\times10^{-6}$, one epoch (403 optimizer
steps) at 7B.

\paragraph{The leak, disclosed up front.} SFT data and every evaluation in
this paper present each board's sixteen words deterministically shuffled
(seeded by puzzle id), precisely so group order is not a giveaway. Our GRPO
loader did not: a one-line bug built training prompts from the answer key,
so during reinforcement learning --- only there --- the words arrived
grouped, positions 1--4 the first group, and so on: every GRPO prompt was
solvable by a positional copy rule worth the full 1.6. Discovered post hoc (\S\ref{sec:control}); disclosed here because nothing in the training-time sections can be read without it. Held-out
measurements are unaffected: evaluation prompts were always shuffled.

\paragraph{The reward.} Table~\ref{tab:reward} is the crux of the paper.

\begin{table}[h]
\centering
\caption{The reward decomposition. Two components are verifiable from the
emitted string alone; three require the answer key. The policy is trained on
the sum.}
\label{tab:reward}
\begin{tabular}{lcl}
\toprule
Component & Value & Cost to check \\
\midrule
Parses as four groups of four & $+0.1$ & \textbf{cheap} --- read the string \\
Correct groups & $+1.0 \times (\text{correct}/4)$ & needs the answer key \\
All four groups correct & $+0.5$ & needs the answer key \\
Group with exactly 3 correct words & $+0.05$ & needs the answer key \\
Invalid output & $-0.1$ & \textbf{cheap} --- read the string \\
\bottomrule
\end{tabular}
\end{table}

Maximum 1.6. The model can raise its score by getting better at the puzzle
or by getting better at emitting well-formed output; only one is the
construct. That asymmetry --- cheap and expensive components, independently
satisfiable --- is the design property this paper studies.

\section{The arc on held-out data}
\label{sec:arc}

Table~\ref{tab:arc} reports all four arms of the reported 7B run on the
162-puzzle test set, greedy decoding, served in a single vLLM session so
every comparison is within-session. The supervised checkpoint is included
because it, not the untrained model, is what GRPO was initialized from.

\begin{table}[h]
\centering
\caption{7B on the test split (162 puzzles, 648 group slots; all four arms
in one serving session; greedy). Means on the 0--4 scale are exact ratios of
the count column. Bootstrap intervals are percentile over puzzles,
$B=1{,}000$ resamples, seed 0. This is the session \S\ref{sec:control}
also scores on the training split.}
\label{tab:arc}
\small
\setlength{\tabcolsep}{4pt}
\begin{tabular}{lccccc}
\toprule
Arm & Groups correct (0--4) & Groups (of 648) & Invalid & Mean reward & Solves \\
\midrule
Instruct (untrained) & 0.1605 $[0.105, 0.222]$ & 26 & 6.8\%  & 0.165 & 0/162 \\
SFT                  & 0.3210 $[0.216, 0.432]$ & 52 & 22.2\% & 0.191 & 2/162 \\
\textbf{GRPO, step 50}  & \textbf{0.4938} $[0.383, 0.623]$ & \textbf{80} & \textbf{16.0\%} & \textbf{0.252} & 2/162 \\
GRPO, step 403     & 0.0247 $[0.000, 0.056]$ & 4 & 0.6\%  & 0.125 & 0/162 \\
\bottomrule
\end{tabular}
\end{table}

Two reference points frame the table. A uniformly drawn valid partition
recovers $4 \times 5775/2{,}627{,}625 = 0.0088$ groups per puzzle, or
\textbf{1.4 slots} across this split; the step-403 endpoint sits at 2.8 times that floor --- and \S\ref{sec:control} shows why: its copied responses score at chance, and the excess over the floor rides on the few responses that are not copies. Second, because all four arms sit
in one session, the stage accounting is paired: $26 \rightarrow 52
\rightarrow 80 \rightarrow 4$, the supervised stage contributing $+0.1605$
$[+0.0556, +0.2716]$ and the reinforcement-learning stage $+0.1728$
$[+0.0741, +0.2840]$ --- near-equal gains, both intervals excluding zero.

The method works, then reverses. Both movements are paired within
Table~\ref{tab:arc}'s session: the climb to step 50 is the two stage gains
just given, each excluding zero, and step 403 falls below the untrained
model by $-0.1358$ $[-0.1975, -0.0802]$. The earlier three-arm session read
$+0.315$ $[+0.204, +0.438]$ and $-0.136$ $[-0.191, -0.080]$ for the same two
contrasts; means in Table~\ref{tab:arc} are exact ratios of the count
column. Two details locate the
mechanism. Held-out mean reward reads 0.252 at step 50 and 0.125 at step
403 against base's 0.165 (a separate session reads step 100 at
0.2664 to a step-50 re-serve's 0.2525),
while the training reward reaches its ceiling by step 125
(\S\ref{sec:mechanism}): the same reward ranks the checkpoints correctly
out of sample and fails to on the training curve --- that contrast, not
the reward's determinism, is this paper's subject. And
the step-50 policy is worse than base on structure (16.0\% invalid against
6.8\%) but \emph{not} worse than its actual initialization, the SFT arm at
22.2\% --- so what is specific to the phase after step 50 is the semantic
collapse, not the structural climb, which was already more than half
complete before the peak (\S\ref{sec:mechanism}).

\paragraph{How wide is the peak?} In a separate serving session, step 100
evaluates at 76/648 against a step-50 re-serve's 80/648, paired difference
$-0.0247$ $[-0.1481, +0.0927]$ ($n=162$). We read this as uninformative
rather than as equivalence: the interval's width is comparable to the
entire step-50 effect, so it is equally consistent with step 100 matching
the peak and with step 100 having lost a third of it, and we did not
pre-specify an equivalence margin. The normalization ablation of
\S\ref{sec:ablation}, this paper's only independent trajectory, argues for
the second reading --- step 100 reads 17/400 against step 50's 35/400 on
the entropy sweep and 35/432 against 62/432 on the greedy checkpoint
sweep, roughly half the peak gone in both. What both sessions agree on is
narrower and is all we claim: the collapse falls between steps 100 and 150,
where the validation semantic score drops 35/400 $\to$ 3/400
(\S\ref{sec:mechanism}).

\paragraph{Reproducibility across serving sessions.} Base and the final
checkpoint reproduce per-puzzle \emph{exactly} across three vLLM sessions
(0/162 differences each --- a deterministic copier and a converged base
leave no argmax ties for batched inference to resolve). Step 50 does not:
2/162 puzzles across sessions (77 vs.\ 80 of 648), 6/162 within one ---
mid-training policies sit near ties, which is why every step-50 number in
this paper carries its serving session. Table~\ref{tab:arc}'s 80/648 is the
four-arm session used throughout; an earlier three-arm session read 77/648,
and the SFT arm decodes to 52 there and 56 in a third
(Appendix~\ref{app:endpoints}). No conclusion in this paper turns on the
difference.

\paragraph{Robustness of the endpoint.} Three independent 7B seeds land at
4, 7 and 11 recovered slots of 648 ($0.045 \pm 0.022$ on the 0--4 scale) ---
all below base's 26. Their LoRA update directions agree pairwise (cosine
$+0.67$ to $+0.69$, merged weight deltas) --- read descriptively, three
runs finding the same attractor, the copy rule (\S\ref{sec:control}); seeds
sharing an initialization and dataset are not isotropic draws, so no
random-direction null applies. Best-of-16 sampling does not rescue the
converged policy: a copier has nothing new to sample (single-seed,
illustrative contrasts in Appendix~\ref{app:endpoints}).

\paragraph{Conditioning on well-formed output widens the gap rather than
explaining it, and we label the operation for what it is.} Is the effect an
artifact of validity, since invalid generations score zero groups and
validity is what GRPO improved? Restricting to
structurally valid generations at 7B, each arm's conditional computed within
a single serving session (the SFT arm's is its 56-slot session): base
\textbf{0.172} ($n=151$ of 162),
SFT \textbf{0.444} ($n=126$), GRPO \textbf{0.025} ($n=161$). This is a
rescaling and we label it as one: each conditional is the arm's
unconditional mean divided by its validity rate, so the gap widens
automatically when the low-scoring arm is the more valid, as here. It
licenses only the negative claim, the one we need: the deficit survives
with formatting held fixed.
Among answers that parse, the converged policy recovers 4 group slots where
the untrained model recovers 26.

\paragraph{The loss shows no stratum gradient we can resolve.} Relative
loss from the untrained model to the endpoint runs 67--100\% across all
four stratum labels (values in Appendix~\ref{app:copy}). We registered
this, pre-diagnosis, as a failed prediction of a selective-inference story;
post-diagnosis it is the copy rule's expected signature, since a positional
rule is blind to stratum (\S\ref{sec:control}). The comparison cannot do
more than that: the strata hold 21--69 test puzzles each, the untrained
model's per-stratum base rates run as low as one recovered slot, and we
report no intervals --- a design with no power to separate uniform loss
from a gradient. Read it as consistent with stratum-blindness, not as
evidence for it.

\section{The scale contrast at 1.5B}
\label{sec:scale}

At 1.5B the instrument reads nothing. The untrained model recovers one
group slot of 648, GRPO one, and SFT --- the best arm --- two, against
Section~\ref{sec:arc}'s 1.4-slot chance floor. No arm is distinguishable
from that floor (SFT's two slots carry an exact binomial interval of
$[0.0015, 0.0444]$, Appendix~\ref{app:small}, which contains it), so no
claim about the construct is available at this scale. What is visible is
the reward driving its verifiable component: invalid output 32.1\% $\to$
2.5\%, mean reward 0.049 $\to$ 0.113 (Appendix~\ref{app:small}). We report
this as an observation: the runs are epoch-asymmetric (806 steps, two
epochs, against 403 and one), we have no 1.5B checkpoint trajectory, and
``the construct was untouched at 1.5B'' would be an endpoint claim,
exactly what this paper shows cannot settle the question.

\section{Trajectory of the collapse: entropy, KL, and the training signal}
\label{sec:mechanism}

Figure~\ref{fig:mechanism} scores checkpoints of the reported run on the
$n=100$ temperature-0.9 \emph{entropy sweep} (denominators of 400; full
per-checkpoint values in Appendix~\ref{app:ckpt}); the ablation's
\emph{checkpoint sweep} (\S\ref{sec:ablation}) is $n=108$, greedy, /432
--- different measurements, never mixed.

\begin{figure}[t]
\centering
\includegraphics[width=\linewidth]{figs/fig1_entropy_collapse.pdf}
\caption{The mechanism, on the $n=100$ temperature-0.9 entropy sweep
(reported 7B run, seed 0). \textbf{Left:} policy entropy against optimizer
step --- a 30-fold collapse, with the cliff between steps 100 and 150.
\textbf{Right:} semantic score against KL from the SFT initialization:
the peak arrives by 2.70 of an eventual 47.05 nats (5.7\%; step 50 is the
earliest checkpoint retained), the remainder undoes it. The untrained
model is an off-trajectory reference at 42.60 nats: KL measures
displacement, not damage (\S\ref{sec:mechanism}).}
\label{fig:mechanism}
\end{figure}

The semantic peak has a location; the structural climb does not. Semantic
rises 0.0325 $\to$ 0.0950 from the SFT initialization to step 50 (13/400
$\to$ 38/400 slots), falls to 0.0075 by step 150, and holds at 0.0100 from
200 on. Structural (valid-partition fraction, not shown) climbs 0.37 $\to$
0.68 $\to$ 0.64 $\to$ 0.95--0.96 across init/50/100/150 on --- one
reversal, at 100 --- with 0.31 of its 0.59 total gain banked by step 50,
concurrent with the semantic peak.

The trajectory is not a trade with a turning point: a broadly monotone
structural climb (one reversal, at 100) with a semantic hump on top, and
only the hump has a location. One gap qualifies the shape: no checkpoint
exists between steps 100 and 150 --- 81\% of the KL displacement (8.12
$\to$ 46.42 nats) --- so the descent is interpolated, not measured.

\paragraph{Entropy collapse is the rule's signature.} Policy entropy falls
13.4-fold between steps 50 and 150, the window in which the semantic score
falls 12.7-fold; by step 150 the policy is effectively deterministic. With
\S\ref{sec:control}'s diagnosis in hand the reading is direct: a positional
copy rule is deterministic, blind to the board's content, and reliably
well-formed: entropy collapse and semantic destruction read as two views of one convergence --- a reading made after the diagnosis, not before it. The structural climb is not --- over half of it ($0.37
\to 0.68$ by step 50) predates measurable copying (0.3\% at step 50 by the strict all-four-quadruples rule of Appendix~\ref{app:copy}; a graded partial-adherence statistic could date copying earlier, which is why we write \emph{measurable}): format gains the rule later locks in. What the collapse
removes is within-prompt variation --- nothing new to sample (the pass@16
result), no group-relative learning signal (the zero-advantage fraction).
And what survives on held-out boards is measured twice over: the cheap
structural term ($0.1(1-2\cdot\text{invalid})$, Table~\ref{tab:arc}) is
52\% / 27\% / 79\% of held-out reward across base / step 50 / endpoint,
and the copy rule itself is visible in 90.7\% of the endpoint's test generations (92.4\% of their emitted groups; \S\ref{sec:control}).

\paragraph{KL is the budget view, and it is a path length, not a damage
scale.} The step-50 checkpoint sits 2.70 nats
from the SFT initialization and is, jointly with step 100
(\S\ref{sec:arc}), the best policy this run produced on
held-out data. The run spent 47.05: 80 slots at 5.7\% of
the displacement, 4 at all of it.
Two qualifications keep this a path description, not a scale: step 50 is
the earliest checkpoint retained, so ``at 5.7\%'' means \emph{at or
before} it; and KL measures displacement, not damage --- the untrained
model sits 42.60 nats out on the same estimator (90.5\% of the run's
displacement) while scoring 0.0125 on this sweep (Appendix~\ref{app:ckpt}),
each estimate under its own policy's samples, so base is off this
trajectory and 5.7\% dates the peak within this run only. The sweep also
compresses base's margin over the endpoint (5/400 vs.\ 4/400 here; 26 vs.\
4 on the greedy test instrument): temperature-0.9 sampling costs a
high-entropy policy far more than a collapsed one.

\begin{figure}[t]
\centering
\includegraphics[width=\linewidth]{figs/fig3_reward_curves.pdf}
\caption{The same reward, in sample and held out (reported 7B run, seed
0). Training reward (gray; sampled $T{=}0.9$, $K{=}8$; TRL 5-step
averages; measured on the answer-ordered prompts of \S\ref{sec:setup})
first touches its 1.6 ceiling at step 125, averaging 1.5814 thereafter.
Held-out mean reward (162 test puzzles; solid markers; greedy; the control session of Table~\ref{tab:control}; right axis, note the scale) peaks at 0.252 at step 50 and ends at 0.125, below the untrained model's 0.165; the hollow step-100 point is a later session (0.2664, vs.\ a step-50 re-serve at 0.2525) and is not joined.}
\label{fig:rewards}
\end{figure}

\paragraph{The training signal, read from the trainer's own logs.} Figure
\ref{fig:rewards} shows the paper's central contrast on one axis: the
training reward first touches the reward function's ceiling of 1.6000 at step
125 and averages 1.5814 over the logged points from then on, with excursions that thin as training proceeds: the deepest is 1.4300 at step 155, the two halves of the 57 post-125 logged points average 1.5663 then 1.5960, and the count of points at exactly 1.6000 rises from 14 of 28 to 24 of 29. We describe this sharpening rather than test it --- the series is autocorrelated and the metric bounded, so no trend statistic we computed is identifiable here, and part of the rise is the instrument (below). What the figure cannot show is TRL's own
\texttt{frac\_reward\_zero\_std} diagnostic: the fraction of prompts whose
eight sampled rewards are identical --- zero group-relative
advantage. It rises from 0.0 at step 50 to 1.0 at step 125
and averages 0.946 over the remainder of training, so from step 125 onward
the optimizer receives no learning signal on nearly every prompt. The saturation is legible under the diagnosis: on answer-ordered prompts (\S\ref{sec:setup}) a copy rule earns 1.6 on every board, seen or unseen --- consistent with saturation beginning mid-epoch, on first-visit batches, and with the vanishing variance. Read as
diagnostics, both in-sample signals bracket the collapse --- the mean's
climb to ceiling (0.9931 $\to$ 1.6000) and the variance's fall to zero both land between logged steps 100 and 125 ---
and neither signs it: saturation with zero spread is what mastery and
degeneracy alike look like. Part of the climb is the instrument: as entropy
collapses, temperature-0.9 sampling converges on greedy, so some of the
rise is the sampling penalty vanishing (the sweep-vs-test gap, reversed). The remaining 278 steps of movement while $\sim$95\% of
prompts carry zero advantage are consistent with continued sharpening of
the rule on the residual prompts that still do; the weight-space
counterpart of that push we have not traced.

\paragraph{$\beta$ was numerically inert, and the bracket is not an
experiment.} An earlier configuration at learning rate $10^{-6}$ and
$\beta=0.04$ froze the policy (KL $\approx 0.001$ for 403 steps; greedy
outputs byte-identical to SFT); the reported one, at $5\times10^{-6}$ and
$\beta=0.001$, overshoots by roughly as much. Two qualifications, both
weakening the pair as an experiment: the configurations differ in
\emph{two} variables at once ($\beta$ by $40\times$, learning rate by
$5\times$), so the pair bounds a region jointly and attributes nothing to
either variable; and $\beta=0.001$ was numerically inert --- $\beta \cdot
\mathrm{KL} = 0.0027$ reward units at step 50, 0.047 at convergence --- so
the KL reference never meaningfully restrained the policy, this paper has
no evidence about a KL penalty large enough to matter, and the
intermediate-$\beta$ sweep remains unrun.

\section{The ablation: the collapse is not the normalization default}
\label{sec:ablation}

We inherited \texttt{scale\_rewards="group"}, verified as the trl 1.10.0
default. Reviewers of an earlier draft independently named it as a
confound: near determinism, the within-group standard deviation dividing
the advantage approaches zero, amplifying sampling noise into large updates
--- which alone might produce the collapse, its timing, and the cross-seed
agreement; \citet{drgrpo} argue against the normalization on independent
grounds.

We reran the 7B pipeline identically --- same SFT warm start, learning
rate, $\beta$, $K$, temperature, LoRA rank, epoch count, seed, and the same
answer-ordered prompts of \S\ref{sec:setup}, so this ablation isolates the
normalization, not the leak --- with
reward scaling disabled, the resolved value read back off the constructed
trainer configuration. The signature is unchanged on the greedy
$n=108$ checkpoint sweep: validation semantic score rises to 62/432 at step
50, falls to 35/432 by step 100 and 4/432 by step 150, while structural
validity climbs 0.70 to 0.96. The validation-selected peak (step 50, frozen
before any test puzzle was scored) recovers 87/648 on test,
$+0.377$ $[+0.247, +0.506]$ paired against base; the final policy 7/648,
below base by $-0.117$ $[-0.179, -0.056]$ paired and within $+0.019$ of
the original endpoint.

The signature survives in the dynamics too, though not point for point:
without reward scaling, policy entropy falls 0.2111 (step 50) $\to$ 0.1346
(100) $\to$ 0.0128 (150) and stays near 0.011 through step 403 (original
endpoint 0.0099), while KL rises 2.75 $\to$ 17.84 $\to$ 45.64 nats per
sequence and plateaus near 46 (original 47.05; 17.84 at step 100 vs.\ the original's
8.12 --- intermediate magnitudes differ). On the $n=100$
temperature-0.9 entropy sweep --- a different measurement from the greedy
$n=108$ checkpoint sweep above --- semantic score reads 35/400 at step 50,
17/400 at 100, then 2/400 from step 150 onward. Hump, collapse, entropy
endpoint, below-base endpoint: all survive with the normalization off. The
phenomenon is not a library default.

\section{The control run: not memorization --- a copy rule}
\label{sec:control}

Training-reward saturation says the policy solves its training prompts
under sampling; was that SFT-stage memorization? The control: all four arms
served in one vLLM session, greedy, scored on the full 807-puzzle training
split and the 162-puzzle test split under the same shuffled prompts as
every other evaluation (Table~\ref{tab:control}).

\begin{table}[t]
\centering
\caption{The control session: four arms, one vLLM session, greedy, both
splits, shuffled evaluation prompts; percentile bootstrap, seed 0. Copy
rate = emitted groups equal to consecutive quadruples of the prompt's word
order (Appendix~\ref{app:copy}).}
\label{tab:control}
\footnotesize
\setlength{\tabcolsep}{4pt}
\begin{tabular}{lccccc}
\toprule
Arm & Train groups (0--4) & of 3{,}228 & Test groups (0--4) & of 648 & Copy rate (test) \\
\midrule
Instruct (untrained) & 0.1871 $[0.156, 0.224]$ & 151 & 0.1605 $[0.105, 0.222]$ & 26 & 2.6\% \\
SFT & 0.5118 $[0.461, 0.568]$ & 413 & 0.3210 $[0.216, 0.432]$ & 52 & 0.3\% \\
GRPO, step 50 & 0.6803 $[0.616, 0.746]$ & 549 & 0.4938 $[0.383, 0.623]$ & 80 & 0.2\% \\
GRPO, step 403 & \textbf{0.0260} $[0.012, 0.041]$ & 21 & 0.0247 $[0.000, 0.056]$ & 4 & \textbf{92.4\%} \\
\bottomrule
\end{tabular}
\end{table}

Three results. The endpoint retains nothing --- 21 of 3{,}228 training
slots, indistinguishable from its own test floor: not a train/test gap, so
not memorization. SFT does partially memorize, though less than the raw
numbers suggest: its $+0.191$ train-minus-test gap (0.5118 vs.\ 0.3210
after three epochs) must be read against the untrained model's own
$+0.027$ gap on the same chronological splits (0.1871 vs.\ 0.1605), which
is split difficulty rather than memorization; the adjusted excess is
$+0.164$. Step 50 inherits essentially all of it ($+0.160$ adjusted) while
improving on SFT by the same margin on both splits ($+0.169$ / $+0.173$):
the peak's gain generalizes. And the shared session yields the
paired stage split: $26 \rightarrow 52 \rightarrow 80 \rightarrow 4$, SFT $+0.1605$ $[+0.0556, +0.2716]$, RL $+0.1728$
$[+0.0741, +0.2840]$, endpoint $-0.1358$ $[-0.1975, -0.0802]$ vs.\ the
untrained model (the original session read $-0.136$ $[-0.191, -0.080]$).

\paragraph{Four lines of string matching find the mechanism.} A 1.6
training reward against 21 recovered training slots forced a look at the
training prompts, where the bug of \S\ref{sec:setup} was found. The behavioral test counts emitted groups
equal, as sets, to consecutive quadruples of the prompt's word order
(definition and full counts in Appendix~\ref{app:copy}). Such quadruples are 91.4\% of the groups the endpoint emits on training boards and 92.4\% on test boards; 88.2\% and 90.7\% of its responses are pure four-quadruple copies; SFT and
step 50 sit at 0.3\%, below even the untrained model's weak bias
(4.9\% / 2.6\%). The converged policy is a prompt-order copier. On training
prompts, where order was the answer, copying earns 1.6; on shuffled prompts
a copy is a fixed valid partition of a random order --- exactly the uniformly drawn valid partition behind \S\ref{sec:arc}'s 1.4-slot floor. Its pure-copy responses score no better than that process: 2 training slots against 6.25 chance-expected --- a three-fold deficit, if anything, and certainly not an excess (we attach no test to it: a puzzle's four slots are not independent, Appendix~\ref{app:ckpt}). What sits above chance is the aggregate --- 21 slots against 7.1 expected; Table~\ref{tab:control}'s puzzle-level interval, $[0.012, 0.041]$ on the 0--4 scale, excludes the chance mean of 0.0088, and test reads 4 against 1.4, same direction, interval not resolving --- and it is carried almost entirely by responses that are \emph{not} pure copies (35 non-copies recover 16 of the 21 slots; 60 partial copies, 3; 711 pure copies, 2): residual solving the rule has not displaced, not copying that works.

\section{The reversal is legible from outside the reward}
\label{sec:judge}

The 162 test generations from base, step 50 and the final policy were
scored blind (board words and completion only) by
three LLM judges (gpt-5-mini, gemini-2.5-flash, claude-sonnet-4-5) on a
1--10 scale, 1458/1458 responses parsed. All three rank step 50 above base
above the final policy, all nine pairwise per-puzzle differences excluding
zero under a paired bootstrap (seed 0; step-50 minus final: $+0.60$
$[+0.44, +0.78]$ / $+3.25$ $[+2.81, +3.67]$ / $+1.62$ $[+1.33, +1.90]$).
That ordering rules out the obvious confound: the final policy is the
\emph{most} well-formed arm on this split (0.6\% invalid against 6.8\% and
16.0\%), so a shared preference for well-formed text would have placed it
first, not last. In hindsight the judges were reading the copy rule: the
arm they unanimously rank last is the copier (\S\ref{sec:control}).
Judges do err in
correlated ways \citep{kohli2026}, so we read this as three correlated
readings of one external signal rather than three independent
confirmations; per-item Pearson $r$ is 0.456 / 0.587 / 0.626 across the
judge pairs (486 items), lowest on the final arm (0.200--0.433); the nine comparisons are one uncorrected family, all excluding zero by margins no standard correction would close. One further limit: the judges score whether the groupings are right --- the reward's semantic component through a different instrument --- so this corroborates the held-out reading rather than adding a second construct.

\section{Implications for evaluation practice}
\label{sec:implications}

\paragraph{Distinguish in-sample from held-out proxy metrics before
concluding a proxy has failed.} Evaluated on held-out data, our reward
ranks the three checkpoints correctly; it fails only in the form
practitioners actually consult it --- as a training curve, in sample, where
it reads perfect. The cheap remedy is not a better reward but the same
reward on data the policy was not updated against --- plus the trainer's
own in-sample brackets (mean saturation, zero variance), though neither
signs the turn (\S\ref{sec:mechanism}). Held out, the separation is real
but not large: paired against the untrained model over the 162 test
puzzles, mean reward moves $+0.086$ $[+0.045, +0.129]$ at step 50 and
$-0.040$ $[-0.061, -0.019]$ at the endpoint, both excluding zero (SFT, at $+0.026$ $[-0.015, +0.068]$, does not). We flag the margin because 79\% of the endpoint's
held-out reward is the cheap structural term: a reward weighted slightly
differently could have kept the aggregate above base while the construct
collapsed exactly as it did. The aggregate caught this run; the
decomposition below is what makes it a diagnostic rather than a
coincidence.

\paragraph{Determinism certifies the scorer, not the task.} Our reward was
correct on every string it scored, and the run was still gamed --- through
the prompt distribution. The audit cost four lines of string matching, but nothing in-sample
prompted it: held-out evaluation said there was something to find.
One function, two measurements: in sample, reliable and invalid --- this
paper's title; held out, what validity it had.

\paragraph{Report reward components separately, never the aggregate.} This
is the recommendation we would keep if we could keep only one. The
failure is invisible in the sum and obvious in the decomposition, which
costs nothing: \S\ref{sec:mechanism} reports it for our three arms (52\%,
27\%, 79\%). It is a standard reward-model diagnostic, simply not applied to
deterministic rewards: determinism is taken to have settled the question.

\paragraph{Treat single-seed RL results as uninterpretable.} The caution
applies to results that come out well too, including step 50: one seed
(endpoint spread across seeds: 0.025--0.068), the argmax of a
nine-point validation grid (step 100 reads 0.0875 to its 0.0950 at
$n=100$, inside one standard error). The grid preceded every test
evaluation of the selected checkpoint --- ordinary validation-based
selection. Whether step 50 is a point or a shoulder we cannot say:
\S\ref{sec:arc}'s step-100 evaluation is too wide to separate the two, and
the ablation trajectory suggests decline has already begun by then.

\section{Related work}
\label{sec:related}

\paragraph{Reward-model overoptimization.} \citet{gao2023scaling}
established this curve shape for \emph{learned} reward models ---
gold-standard performance rising then declining with KL --- where the
explanation is Goodhart against an approximation \citep{stiennon2020}. Our reward is a program: it approximates nothing, and
there was nothing in the \emph{function} to game --- the gaming ran through
the prompt distribution (\S\ref{sec:setup}), a channel the learned-reward
literature's remedy, a better reward model, does not touch. Determinism is
the very feature cited to argue verifiable rewards escape this failure
mode; here it guaranteed only that the shortcut was scored accurately.

\paragraph{RL for reasoning and entropy collapse.} Verifiable-reward RL
underlies recent reasoning systems \citep{deepseekr1, deepseekmath};
\citet{yue2025} ask whether it extends capability beyond the base model's pass@$k$ ceiling, and here RLVR, initialized from an instruction-tuned stack, ends below \emph{both} of its own antecedents on the construct. \citet{shao2025spurious} show Qwen2.5 models gaining on MATH-500 under RLVR with random or incorrect rewards, by amplifying latent base-model behaviors; part of our step-50 gain, obtained under a corrupted signal, may be that phenomenon --- the peak is not evidence the signal was sound.
\citet{cui2025} document entropy collapse as a central dynamic of RL for
reasoning; here the collapse coincides with convergence onto a leaked shortcut, and we read it as symptom rather than independent cause --- a post-hoc reading the shuffled-prompt rerun can test (\S\ref{sec:control}, \S\ref{sec:limitations}). \citet{drgrpo} critique group-level advantage normalization;
\S\ref{sec:ablation} shows the phenomenon does not depend on it.

\paragraph{Reward hacking and its taxonomy.} \citet{skalse2022}
formalize the failure mode (as \emph{reward gaming}); ours is gaming in
their sense against the composed task --- reward plus prompt distribution
--- not the reward function, which was exact.
\citet{abahana2026} classifies RLHF failure transitions by directional
changes in learned reward and judge scores; ours differs in instrument (a
rule-based reward, exact) and in detection (bracketed by held-out data
rather than flagged in flight). \citet{krakovna2020specification} catalogue
specification gaming across RL systems; the entry our run would occupy is
an ordinary one, and we claim novelty for the trace rather than the
phenomenon.

\paragraph{Shortcut learning and leakage.} What our policy learned is a
shortcut in the sense of \citet{geirhos2020shortcut}: a decision rule that
succeeds on the training distribution and fails off it, because a feature
of the input --- here, word order --- carried the label. \citet{lapuschkin2019unmasking}
make the same point for supervised vision models under the name Clever Hans
prediction. Neither literature is usually applied to RL post-training, where the training distribution is regenerated by a trainer, so a leak hides in code rather than in data. That is the transfer we are making.

\paragraph{Abruptness.} \citet{pan2022effects} report phase transitions in reward misspecification --- capability thresholds at which true reward falls sharply --- in some environments they study. Our collapse between steps 100 and 150 has that shape, which is why the gap in our checkpoint grid is a real limitation (\S\ref{sec:limitations}), not an interpolation detail: if the turn is a transition, its location is the measurement worth having.

\paragraph{Validity in AI evaluation.} Measurement-theoretic critiques of
ML evaluation \citep{jacobs2021, freiesleben2025} distinguish a measurement
from the construct it operationalizes; \citet{chehbouni2025} apply it to LLM
judges; ours is a case where the measurement is maximally reliable and the
validity gap does the damage anyway. \citet{agarwal2021} motivate our
statistical practice.

\section{Limitations}
\label{sec:limitations}

The largest limitation: no leak-free run exists. Every training run here
(the reported 7B, the three seeds, the ablation) consumed the
answer-ordered prompts, so nothing separates what GRPO does to this
construct from what it did to the leaked presentation of it; the one-line fix is in the repository, and the shuffled-prompt
rerun is the experiment these conclusions wait on. One task, one reward design; the
checkpoint trajectory is 7B seed 0 only, so the scale claim rests on
endpoints at both scales and a trajectory at one, and an endpoint cannot
rule out a traversed-and-returned 1.5B peak (\S\ref{sec:scale}). No human baseline exists: ``worse than the untrained model'' is
established, ``bad in absolute terms'' is not. Three-seed evidence covers only the converged endpoint (spread
0.025--0.068, why we distrust single numbers); the peak's height and
location may move across seeds, and our step-100 evaluation is too wide to
establish that it is not a knife edge (\S\ref{sec:arc}). A single $\beta$ is reported alongside one that froze the policy (no
intermediate run, a learning-rate difference besides, $\beta \cdot
\mathrm{KL} \le 0.047$ throughout), so a working KL penalty remains
untested. The diagnosis is post hoc
(\S\ref{sec:mechanism}'s readings are retrospective), the 1.5B loader
shares the leak (copy rate unmeasured), and the blind judges are
correlated scorers, not independent votes (\S\ref{sec:judge}).


\bibliographystyle{plainnat}
\bibliography{refs}

\appendix

\section{The two prompts}
\label{app:prompts}

The bug of \S\ref{sec:setup} is one line, and it is easiest to see as the
pair of strings the model actually received. Both are the user turn of the
same puzzle (id 3, 2023-06-14), under the same system prompt; the answer key
is \texttt{\{CHEEK, EYE, MOUTH, NOSE\}}, \texttt{\{CHOW, GOBBLE, SCARF,
WOLF\}}, \texttt{\{LAB, PEKE, PIT, POM\}}, \texttt{\{AMIGO, KING, STOOGE,
TENOR\}}.

\medskip
\noindent\textbf{GRPO training prompt, as the reported run consumed it}
(\texttt{train/grpo.py::build\_dataset}):

{\footnotesize\begin{verbatim}
Words: CHEEK, EYE, MOUTH, NOSE, CHOW, GOBBLE, SCARF, WOLF,
       LAB, PEKE, PIT, POM, AMIGO, KING, STOOGE, TENOR
\end{verbatim}}

\noindent\textbf{SFT data and every evaluation}
(\texttt{data/formatting.py::shuffled\_words}, a puzzle-seeded permutation):

{\footnotesize\begin{verbatim}
Words: CHOW, POM, PIT, KING, AMIGO, NOSE, SCARF, CHEEK,
       EYE, TENOR, STOOGE, GOBBLE, MOUTH, LAB, PEKE, WOLF
\end{verbatim}}

\medskip
In the first, emitting the words four at a time in the order given scores
the full 1.6. In the second it scores what a uniformly drawn valid
partition scores. Nothing else differs --- same system prompt, same reward,
same parser --- and the difference is a missing call to the shuffle the
rest of the codebase already used. The regression test we have since added
asserts the two paths agree per puzzle, which is the property that would
have caught this and the weaker ``is not in answer order'' assertion would
not.

\section{Per-checkpoint values for the entropy sweep}
\label{app:ckpt}

Table~\ref{tab:ckpt} tabulates all eleven rows of the sweep artifact: the
trajectory plotted in Figure~\ref{fig:mechanism} (including the step-400
checkpoint, indistinguishable from step 403 at this precision), the
structural column discussed in Section~\ref{sec:mechanism}, and the
off-trajectory untrained Instruct reference. All values are from the $n=100$
temperature-0.9
entropy sweep of the reported 7B run (seed 0); the artifact is
\texttt{results-analysis/entropy-kl-7b.json} in the released repository,
\repourl{}, which also holds the
training and evaluation configs, the reward implementation, and the
generation logs behind every table here.

Unless stated otherwise, all intervals in this paper are percentile
bootstrap over \emph{puzzles} --- the unit of independent sampling --- with
$B=1{,}000$ resamples at seed 0; per-slot counts are reported alongside
means but never resampled independently, since a puzzle's four slots are
not independent. Where a count is small enough that the bootstrap
degenerates (Table~\ref{tab:endpoints}), we substitute Clopper--Pearson
intervals and say so.

\begin{table}[h]
\centering
\caption{The reported 7B run on the $n=100$ temperature-0.9 entropy sweep.
Structural = fraction of outputs that are a valid partition; semantic = mean
correct groups / 4. KL is the per-sequence sample-based estimate against the
SFT initialization, each row estimated under its own policy's samples; the
untrained Instruct row is off the GRPO trajectory
(\S\ref{sec:mechanism}).}
\label{tab:ckpt}
\begin{tabular}{lcccc}
\toprule
Checkpoint & Policy entropy & KL from SFT init & Structural & Semantic \\
\midrule
Instruct (untrained) & 0.2577 & 42.60 & 0.69 & 0.0125 \\
\midrule
SFT init  & 0.3025 & 0.00  & 0.37 & 0.0325 \\
Step 50   & 0.2107 & 2.70  & 0.68 & \textbf{0.0950} \\
Step 100  & 0.1999 & 8.12  & 0.64 & 0.0875 \\
Step 150  & 0.0158 & 46.42 & 0.95 & 0.0075 \\
Step 200  & 0.0109 & 46.88 & 0.95 & 0.0100 \\
Step 250  & 0.0102 & 46.96 & 0.96 & 0.0100 \\
Step 300  & 0.0099 & 47.05 & 0.96 & 0.0100 \\
Step 350  & 0.0099 & 47.05 & 0.96 & 0.0100 \\
Step 400  & 0.0099 & 47.05 & 0.96 & 0.0100 \\
Step 403  & 0.0099 & 47.05 & 0.96 & 0.0100 \\
\bottomrule
\end{tabular}
\end{table}

\section{7B endpoints across three seeds}
\label{app:endpoints}

Table~\ref{tab:endpoints} tabulates the three-seed endpoint session
discussed in Section~\ref{sec:arc}; the 4 / 7 / 11 spread is the basis of
the single-seed caveats in Sections~\ref{sec:implications}
and~\ref{sec:limitations}. The pass@16 contrasts referenced there are
differences of $+0.920$ $[0.778, 1.062]$ for SFT over GRPO and $-0.370$
$[-0.475, -0.259]$ for GRPO against the untrained Instruct model ---
single-seed, read as illustrative only.

\begin{table}[h]
\centering
\caption{7B endpoints across three seeds (separate serving session from
Table~\ref{tab:arc}; the SFT arm decodes to 0.346/56 in the
Table~\ref{tab:arc} session's configuration and 0.321/52 here, differing on
2 of 162 puzzles; the untrained and final-checkpoint anchor arms reproduce
across sessions exactly).}
\label{tab:endpoints}
\small
\setlength{\tabcolsep}{4pt}
\begin{tabular}{lcccc}
\toprule
Arm & Groups correct (0--4) & Groups (of 648) & Invalid & Mean reward \\
\midrule
Instruct (untrained) & 0.160 & 26 & 6.8\% $[3.1, 11.1]$ & 0.165 \\
SFT   & 0.321 & 52 & 22.2\% & 0.197 \\
GRPO (3 seeds) & $0.045 \pm 0.022$ & 4 / 7 / 11 & $1.2 \pm 0.6$\% & $0.132 \pm 0.008$ \\
\bottomrule
\end{tabular}
\end{table}

\section{Control-session copy rates and rewards}
\label{app:copy}

Table~\ref{tab:copy} gives the full copy-rule test behind
Section~\ref{sec:control}: for each generation, an emitted group counts as
a copy when it equals, as a set, a consecutive quadruple (positions 1--4,
5--8, 9--12, 13--16) of that prompt's word order. Train prompts parsed:
806 of 807 per arm; test: 162 of 162. Mean held-out rewards in this
session: base 0.1802 (train) / 0.1654 (test), SFT 0.2490 / 0.1910, step 50
0.3113 / 0.2519, endpoint 0.1242 / 0.1250 --- the endpoint's test values
reproduce the published 0.125 exactly. The per-stratum relative loss cited
in \S\ref{sec:arc} (untrained model $\to$ endpoint, over the 158 of 162
test puzzles with a single stratum label): 92\% on \texttt{category}
($n=69$), 67\% on \texttt{cultural} ($n=28$), 86\% on \texttt{tag-fillin}
($n=40$), 100\% on \texttt{wordplay} ($n=21$, from a near-floor base of
0.048, one recovered slot).

\begin{table}[h]
\centering
\caption{Copy-rule rates, control session (greedy; shuffled evaluation
prompts). Group-level = share of emitted groups that are prompt-order
quadruples; response-level = share of responses whose four groups all are.}
\label{tab:copy}
\small
\setlength{\tabcolsep}{4pt}
\begin{tabular}{lcccc}
\toprule
Arm & Train (group) & Train (response) & Test (group) & Test (response) \\
\midrule
Instruct (untrained) & 4.9\% (159/3224) & 4.3\% & 2.6\% & 1.2\% \\
SFT & 0.3\% (9/3224) & 0.0\% & 0.3\% & 0.0\% \\
GRPO, step 50 & 0.3\% (11/3224) & 0.0\% & 0.2\% & 0.0\% \\
GRPO, step 403 & \textbf{91.4\%} (2946/3224) & \textbf{88.2\%} & \textbf{92.4\%} & \textbf{90.7\%} \\
\bottomrule
\end{tabular}
\end{table}

\section{1.5B test-split table}
\label{app:small}

Table~\ref{tab:small} tabulates the 1.5B arms discussed in
Section~\ref{sec:scale}; every mean is the exact ratio of its count to 162.

\begin{table}[h]
\centering
\caption{1.5B on the test split (162 puzzles, 648 group slots).
Groups-correct intervals are Clopper--Pearson exact binomial on the slot
count of 648, scaled to 0--4 (a percentile bootstrap is degenerate at counts
of 1--2); slots cluster within puzzles, so they are approximate.
Invalid-rate intervals are percentile bootstrap over puzzles. Conditioned on
structurally valid output (\S\ref{sec:arc}), the arms read: base 0.009
($n=110$ of 162), SFT 0.048 ($n=42$), GRPO 0.006 ($n=158$).}
\label{tab:small}
\small
\setlength{\tabcolsep}{4pt}
\begin{tabular}{lccccc}
\toprule
Arm & Groups correct (0--4) & Groups (of 648) & Invalid & Mean reward & Solves \\
\midrule
Base & 0.0062 $[0.0002, 0.0343]$ & 1 & 32.1\% $[24.7, 38.9]$ & 0.049 & 0.0\% \\
SFT  & 0.0123 $[0.0015, 0.0444]$ & 2 & 74.1\% $[67.3, 80.2]$ & $-0.038$ & 0.0\% \\
GRPO & 0.0062 $[0.0002, 0.0343]$ & 1 & 2.5\% $[0.6, 5.6]$ & 0.113 & 0.0\% \\
\bottomrule
\end{tabular}
\end{table}

\end{document}
